\section{问题分析}

\subsection{问题一分析}

问题一由需求统计、时间序列预测和末端基础调度三个环节构成。首先，应核验任务字段、时间单位、区域编码及任务类型映射，并按区域—任务类型聚合任务数、GPU 需求、GPU-hour、执行时长、到达强度和并发水平，以识别趋势、周期、波动和区域相关性。由于任务序列具有严格时间顺序，训练、验证和测试必须按时间切分，标准化、特征筛选及超参数选择均不得读取测试区间信息。预测目标被统一为逐小时 GPU 需求，并通过 WAPE、MAE 等指标进行评价；当真实需求和为零或样本方差不足时，相应比例误差或相关性不应被强行计算。

预测模型确定后，第 $2376$ 至 $2399$ 小时的正式调度仍应读取同期实际到达任务，而不能以预测结果替代真实任务。对于持续时间并非整数小时的任务，需要计算任务执行区间与每个小时的交集，由此形成 GPU-hour 和 AI IT 负荷。实时任务的开工时刻被固定为到达时刻，弹性任务则在到达时刻、截止时点和系统终点共同限定的完整整数时域内寻找可行位置。调度过程还需同步检查有向网络时延、GPU 容量、总 IT 功率及设施功率，从而保证甘特图中的实际起止时刻与逐时资源利用率保持一致。

\subsection{问题二分析}

问题二的输入为第 $0$ 至 $2399$ 小时的实际任务以及第 $0$ 至 $2406$ 小时的电力参数，输出为全部任务的执行区域、开工与完成时刻及对应能源指标。其关键在于任务位置变化会同时改变区域 GPU 占用、AI IT 功率、PUE 损耗、新能源直接消纳和电网购电，因此不能仅依据低碳强度或高新能源出力进行孤立迁移。区域总 IT 负荷应由可迁移 AI 负荷与固定非 AI 负荷相加得到，设施负荷则由总 IT 负荷乘 PUE 形成。

该问题规模较大，且任务不可抢占约束造成跨小时连续占用。为获得可执行方案，先依据实时优先、截止时间和持续时间等规则构造全时域可行初始解，再对高 GPU-hour 弹性任务进行有界确定性局部改进。候选位置应由成本、碳排放、剩余新能源、等待时间和网络时延的无量纲评分排序，并在每次移动前移除旧占用、确认新位置可行后再更新系统状态。最终成本、碳排放、GPU-hour 加权网络时延及新能源利用率均应从完整候选重新计算，所得方案只能解释为有限搜索范围内的情景可行候选。

\subsection{问题三分析}

问题三将任务调度与逐时 IT 负荷视为已知条件，因而优化对象被限定为新能源直接供能、新能源充电、电网充电、储能放电、购电、外送和弃电等能源流。设施负荷应由基线 AI IT 负荷与非 AI IT 负荷之和乘区域 PUE 得到，不能调用问题二调整后的任务负荷。储能状态必须从第 $0$ 小时运行前的已知初始 SOC 开始逐时递推，附件中的基准 SOC 轨迹只用于比较，不应替代优化状态。

储能策略需要兼顾价格套利、新能源消纳、削峰和减小负荷波动，但这些目标可能相互冲突。实际计算中可依据各区域完整购电价格序列的四分位数识别低价充电与高价放电时段，并按固定能源流顺序生成确定性可行候选。每一时段均需核验新能源守恒、设施负荷平衡、充放电互斥、购售电边界及 SOC 上下限，同时限制售电来源，避免电网购电转售。峰值净购电和负荷波动应由最终净购电序列复算，不能由局部能源流直接推断。

\subsection{问题四分析}

问题四将前三问中的任务约束、能源平衡和储能动态合并为统一的算—储—电协同问题。其输入仍为实际任务、实际逐时电力参数和统一 GPU 功率映射，输出包括任务迁移与开工策略、储能轨迹、购售电及新能源分配方案，以及成本、碳排放、网络时延、服务质量、新能源利用率和区域峰值净购电等评价指标。服务质量只能由即时开工、SLA 时延、按期完成和弹性任务等待时间等可观测量构造，因此只能作为代理指标解释。

任务二元决策与储能连续状态横跨完整时域，直接穷举全部组合不可行。因而，任务侧采用保持全时域可行性的有界局部改进，能源侧采用从初始 SOC 前向递推的确定性储能规则，并从同一完整候选统一复算全部目标。情景比较中，仅改变碳约束、电价机制或新能源出力，其他任务集、设备容量、网络时延、功率映射和评价口径保持不变。不同情景所得候选需要经过有限值、资源容量、能量平衡和终端状态复核；若声明范围内未找到满足目标的候选，则应报告该目标在既定启发式范围内不可达。

\subsection{子问题间关系分析}

问题一的统计和预测为识别任务规律提供依据，其任务重叠计算和可行调度框架可被后续问题复用，但预测结果不直接进入问题二至问题四的正式调度。问题二在实际任务基础上扩展碳感知迁移与开工优化，问题三则因负荷已由附件固定，可与前两问相对独立地开展。问题四继承任务侧和能源侧约束，但必须重新形成统一候选并复算指标，不能将问题二的任务方案与问题三的储能方案直接拼接。
